\subsection{Sprint 3}

\subsubsection*{Planejamento e Distribuição das Atividades}

Com a definição das tarefas previstas para a Sprint 3, a Squad Ninja ficou responsável
pela implementação da funcionalidade de \textbf{Gestão de Currículos}, destinada ao
uso administrativo da plataforma.

A divisão das responsabilidades ocorreu da seguinte forma: Mayra (\ac{ages} I) ficou
responsável pelo desenvolvimento do backend da funcionalidade; Vinicius (\ac{ages} III)
atuou no acompanhamento das atividades da equipe, conferência de \textit{Pull Requests}
e suporte para esclarecimento de dúvidas; enquanto eu fiquei responsável pelo
desenvolvimento do frontend da task. Além disso, também assumi uma atividade
relacionada a débito técnico, denominada \textbf{Desmock do Dashboard da Administração}.

\subsubsection*{Desenvolvimento da Funcionalidade de Gestão de Currículos}

No dia seguinte ao planejamento da sprint, iniciei o desenvolvimento da funcionalidade
a partir da descrição presente no épico:

\begin{quotation}
\itshape
``Será adicionado um botão, na coluna ações, na tabela de gestão de alunos. Ao clicar,
o administrador será redirecionado para visualização de currículo daquele estudante.''
\end{quotation}

Com base nessa descrição, interpretei que seria necessário implementar um botão para
visualização dos currículos dos alunos. Considerando o contexto da funcionalidade,
desenvolvi a solução tanto para a tela de Gestão de Alunos quanto para a tela de Gestão
de Currículos, realizando as alterações necessárias no frontend e no backend.

Após concluir a implementação, informei ao \ac{ages} IV responsável pela atividade.
Posteriormente, foi esclarecido que havia ocorrido uma inconsistência na comunicação
dos requisitos e que o botão deveria estar presente apenas na tela de Gestão de
Currículos. Mesmo assim, durante as discussões realizadas, foi considerado que a
funcionalidade poderia permanecer implementada localmente para futuras avaliações,
uma vez que não impactava negativamente a interface e poderia representar uma melhoria
de usabilidade.

Durante esse processo também foram identificados alguns problemas relacionados aos
filtros da aplicação, que foram reportados à equipe e posteriormente corrigidos.

\subsubsection*{Implementação da Tela de Gestão de Currículos}

Na semana seguinte, foi desenvolvido o frontend da tela de Gestão de Currículos. Para
acelerar a implementação e manter a consistência visual do sistema, utilizou-se como
referência a tela já existente de Gestão de Alunos, uma vez que ambas possuíam
estrutura e campos semelhantes.

Durante o desenvolvimento foi identificada uma diferença relacionada ao componente de
paginação, que ainda utilizava um padrão anterior do projeto. Essa situação ocorreu
porque a branch local não estava completamente atualizada com a versão mais recente
da branch de desenvolvimento. Entretanto, essa diferença não comprometeu o andamento
da implementação nem a integração da funcionalidade.

\subsubsection*{Desenvolvimento da Task de Débito Técnico}

Durante a análise da task de débito técnico, identifiquei uma possível semelhança entre
a atividade que me havia sido atribuída e outra que já estava em andamento por outro
integrante da equipe. Diante dessa situação, busquei alinhamento com os envolvidos
para confirmar se ambas as demandas eram realmente distintas.

Na ocasião, o entendimento do meu colega era de que as tarefas possuíam escopos
diferentes, principalmente devido às diferenças observadas nas descrições registradas.
Com base nessas informações, prossegui com a implementação da funcionalidade e
realizei a abertura do \textit{Pull Request} para revisão.

Posteriormente, durante o processo de conferência, foi constatado que ambas as
atividades atendiam à mesma necessidade do sistema, resultando em uma implementação
duplicada. Como consequência, o \textit{Pull Request} foi encerrado sem integração ao
projeto.

Essa situação reforçou a importância de uma comunicação mais centralizada entre os
responsáveis pela definição das tarefas e os desenvolvedores, bem como da manutenção
de descrições mais detalhadas e alinhadas entre as diferentes fontes de planejamento.

\subsubsection*{Avaliação da Sprint}

De maneira geral, a Sprint 3 apresentou resultados bastante positivos. Foi possível
perceber uma evolução significativa na maturidade das squads e na capacidade de entrega
da equipe. Comparativamente às sprints anteriores, houve um aumento na quantidade de
tarefas desenvolvidas dentro do mesmo período de tempo, sem prejuízo à qualidade das
entregas.

Outro aspecto que demonstrou evolução foi a documentação das atividades no ClickUp.
As tarefas passaram a apresentar descrições mais completas e detalhadas, reduzindo
ambiguidades e facilitando a compreensão dos requisitos. Em muitos casos, além dos
épicos principais, as subtarefas continham detalhamentos adicionais que auxiliavam
diretamente no desenvolvimento.

Por fim, a apresentação das funcionalidades para a stakeholder demonstrou que as
entregas realizadas estavam alinhadas às expectativas do projeto. O retorno recebido
evidenciou satisfação com os resultados alcançados e reforçou a percepção de que as
demandas apresentadas estavam sendo atendidas de forma consistente pela equipe.
